%%%%%%%%%%%%%%%%%%%%%%%%%%%%%%%%%%%%%%%
% Trabalho de Sistemas Distribu�dos
%%%%%%%%%%%%%%%%%%%%%%%%%%%%%%%%%%%%%%%

\documentclass[11pt,brazil, a4paper]{article}

%%%%%%%%%%%%%%%%%%%%%%%%%%%%%%%%%%%%%%%
% Preambulo - Configura��o de p�gina
%%%%%%%%%%%%%%%%%%%%%%%%%%%%%%%%%%%%%%%
\usepackage{color}
\usepackage[portuges]{babel}
\usepackage[T1]{fontenc}
\usepackage[ansinew]{inputenc}
\usepackage{array}
\usepackage{colortbl}
\usepackage{epsfig}
\usepackage{multirow}
\usepackage{hyperref}
\usepackage{amsmath}

\usepackage{tikz}
\usetikzlibrary{3d}
\usetikzlibrary{calc,decorations.pathmorphing} % noisy shapes
\usetikzlibrary{fit}					% fitting shapes to coordinates
\usetikzlibrary{backgrounds}	% drawing the background after the foreground
\usetikzlibrary{shadows}
\usetikzlibrary{mindmap}
\usetikzlibrary{positioning,shapes,shadows,arrows}
\usepackage{pgfplots}
\pgfplotsset{compat=1.8}
\usetikzlibrary{arrows,shapes}

\pagestyle{empty}

% Define margem no topo como sendo 1,7 cm
\setlength{\topmargin}{0pt}
\setlength{\voffset}{-1.8cm}

%define margem inferior como 1,5 cm
\setlength{\textheight}{25,6cm}

%define margens laterais como sendo 1,5 cada
\setlength{\oddsidemargin}{0cm}
\setlength{\hoffset}{-1.2cm}
\setlength{\textwidth}{18cm}

%Define campo com sombreamento cinza
\newcolumntype{G}{>{\columncolor[gray]{0.8}[\tabcolsep]\tiny \bf \sf}p{13pt}<{\centering}}
\newcolumntype{T}{>{\tiny \sf}p{13pt}<{\centering}}


\usepackage{xcolor}
\definecolor{verde}{rgb}{0,0.5,0}

\usepackage{listings}
\lstset{
	language=c,
	basicstyle=\ttfamily\normalsize, 
	keywordstyle=\color{blue}, 
	stringstyle=\color{verde}, 
	commentstyle=\color{red}, 
	extendedchars=true, 
	showspaces=false, 
	showstringspaces=false, 
	numbers=left,
	numberstyle=\tiny,
	breaklines=true, 
	breakautoindent=true, 
	captionpos=b,
	xleftmargin=0pt,
}


%%%%%%%%%%%%%%%%%%%%%%%%%%%%%%%%%%%%%%%
% Inicio do documento
%%%%%%%%%%%%%%%%%%%%%%%%%%%%%%%%%%%%%%%
\begin{document}
	
	\selectlanguage{brazil}
	
	\hspace{-0.8cm} \begin{tabular*}{18cm}{p{2cm}p{12cm}@{\extracolsep{\fill}}}
		\multirow{2}{2cm}{\epsfig{file=puclogo_small_bw,width=1.5cm}}
		&  \vspace{1mm} \usefont{T1}{cmss}{bx}{n} {\centering PONTIF�CIA UNIVERSIDADE CAT�LICA DE MINAS GERAIS }\\
		&  \usefont{T1}{cmss}{x}{n} Instituto de Ci�ncias Exatas e Inform�tica \vspace{3mm}  \\
	\end{tabular*}
	
	
	%%%%%%%%%%%%%%%%%%%%%%%%%%%%%%%%%%%%%%%
	% Define caracteristicas da disciplina
	%%%%%%%%%%%%%%%%%%%%%%%%%%%%%%%%%%%%%%%
	\hspace{-0.8cm} \begin{tabular*}{18.5cm}{|p{7cm}@{\extracolsep{\fill}}|p{6cm}|p{1.7cm}|}
		\hline
		\sf Disciplina & \sf Curso  & \sf Per�odo \\
		
		\multicolumn{1}{|r|}{Algoritmos e Estruturas de Dados II} &
		\multicolumn{1}{c|}{Engenharia de Computa��o} &
		\multicolumn{1}{c|}{2$^\circ$} \\ \hline \hline
		\multicolumn{3}{|l|}{\sf Professor} \\
		\multicolumn{3}{|l|}{\hspace{1cm}Kleber Jacques F. de Souza (klebersouza@pucminas.br)} \\ \hline
	\end{tabular*}
	
	
	\vspace{0.1cm}
	
	
	\begin{center}
		{\Large \bf Lista de Exerc�cios - Hashing e �rvore Bin�ria Balanceada} \\
		%{\normalsize \bf em dupla}
	\end{center}
	

%===================================================================================================


\begin{enumerate}

\item Implemente uma estrutura de pesquisa usando a Tabela Hashing, usando as seguintes defini��es:

\begin{enumerate}
	\item Tratamento de colis�o usando Endere�amento Aberto, onde os registros em conflito s�o armazenados dentro da pr�pria tabela, executando uma busca sequencial.
	\item Tratamento de colis�o usando Endere�amento Separado, onde os registros em conflito s�o armazenados dentro de uma lista encadeada.
\end{enumerate}

\item Crie uma �rvore Bin�ria AVL com a inser��o dos seguintes itens, mostrando a constru��o da �rvore passo a passo:\\

5 - 10 - 15 - 30 - 40 - 20 - 25 - 35 - 45 - 50

\item Considere que voc� tem os seguintes itens a serem armazenados no sistema: \\

9 - 8 - 7 - 6 - 5 - 4 - 3 - 2 - 1 - 0 \\

Fa�a a modelagem para armazenar essas informa��es, na ordem em que aparecem, nas seguintes estruturas de dados:

\begin{itemize}
	\item Pesquisa Sequencial
	\item Pesquisa Bin�ria
	\item �rvore Bin�ria
	\item �rvore Bin�ria Balanceada AVL
	\item Tabela Hash
\end{itemize}

Usando a modelagem das estrutura de pesquisa, calcule o tempo de acesso para pesquisar os seguintes valores, em cada uma das estruturas:
\begin{enumerate}
	\item 9
	\item 7
	\item 5
	\item 3
	\item 1
\end{enumerate}


\end{enumerate}






%===================================================================================================

\end{document}